% 1. Discuss functional concepts
% 2. Mention connection to symbols
% 3. Cite Pfleger and Stacuzzi incremental work
% 4. Discuss relation to schemas, frames, skills, others
% 5. Discuss chunks in Soar and ACT-R 

\documentclass[11pt,letterpaper]{article}
\usepackage{cogsys}
% \usepackage{apacite}
% \usepackage{graphicx}
\usepackage[T1]{fontenc}
\usepackage{times}
\usepackage[pdftex]{graphicx} % use this when importing PDF files
\usepackage{epsf}

% Various personal itemization commands. 
\def\texitem#1{\par\noindent\hangindent 20pt
               \hbox to 20pt {\hss #1 ~}\ignorespaces}
\def\itemitem#1{\par\noindent\hangindent 40pt
                \hbox to 40pt {\hss #1 ~}\ignorespaces}
\def\bull{\texitem{$\bullet$}}
\def\bbull{\itemitem{$\circ$}}
\def\note#1{\vskip 0.15in\texitem{$\bullet$} #1 \vskip 0.15in}

\def\prodigy/{{\sc Prodigy}}
\def\icarus/{{\sc Icarus}}

\cogsysheading{11}{2025}{1--12}{8/2025}{10/2025}

\ShortHeadings{Concepts and Chunks in Cognitive Systems} 
              {P.\ Langley}

\begin{document} 

\title{Concepts and Chunks in Cognitive Systems} 
 
\author{Pat Langley}{pat.langley@gtri.gatech.edu}
\address{Information and Communications Laboratory, Georgia Tech 
         Research Institute, Atlanta, GA 30318} 
\vskip -0.13in
\address{Institute for the Study of Learning and Expertise,
         2164 Staunton Court, Palo Alto, CA 94306}
\vskip 0.2in

\begin{abstract}
Research in cognitive science, both empirical and theoretical, has
implicitly distinguished between two types of mental structures:
concepts and chunks. The first paradigm has studied people's ability
to represent, recognize, and learn categories, with classification error
and response time as typical measures. The second has examined humans' 
ability to store, use, and acquire recurring patterns, with memory
capacity as the main dependent variable. The literatures on these two
phenomena are nearly disjoint, but they also share many assumptions.
In this essay, I review the similarities and differences between
computational accounts of concepts and chunks, then propose some
steps toward developing a unified theory of their roles in natural
and artificial cognitive systems.
\end{abstract}  

\vspace{0.1in}
\section{Introduction} 

% Review cogsys paradigm, emphasis on cog structures
% Examples of structures (HLL paper, PUG)
% Structures not enough, also need organization, performance, 
% learning (give egs) 

The cognitive systems paradigm has a number of characteristics that
distinguish it from other computational approaches to intelligence
(Langley, 2012). Perhaps the most basic feature is its emphasis on
the central role of {\it cognitive structures\/}. Of course, these
alone do not suffice to explain mental abilities; they must also be
organized in some manner, performance mechanisms must use them 
effectively, and learning processes must acquire and refine them. 
However, these all build on cognitive structures, so discussions
naturally start with them. 

% Aim for general theories of mind (e.g., cog archs)
% These often involve multiple structure types but no more than needed
% Thus, we should look for opportunities to unify ideas

Like other scientific disciplines, the cognitive systems movement aims 
for general theories with broad coverage of phenomena, in this case
abilities found in human intelligence. Such coverage appears to require
multiple types of mental structures and mechanisms that operate on them,
as illustrated in many cognitive architectures (Langley, Laird, \&
Rogers, 2009). At the same time, we should not postulate different types 
of content unless necessary, as elegant theories are more desirable than 
complicated ones, other things being equal. 

% In this essay, will discuss prospects for unifying two: concepts and chunks
% First review each, including rep, perf, learning, and abilities afforded
% Next discuss differences, similarities, and potential for unification
% Finally, outline 

In this essay, I discuss the prospects for unifying two classic types of
cognitive structures -- {\it concepts\/} and {\it chunks\/} -- that have
typically been treated separately by cognitive psychologists and AI
researchers. I review computational accounts of these two notions in
turn, including common assumptions about representation, organization,
performance, and learning. Next I summarize similarities and differences
in these narratives, after which I propose an agenda for developing unified
theories that would cover the phenomena associated with both types of
structures. In closing, I outline three promising alternatives for
tackling this unification challenge. 

% \vspace{-0.10in}
\section{Concepts in Cognitive Systems}

% Definition

The study of concepts and categories has been a mainstay of research
on cognitive psychology since the 1960s, as discussed by Rips, Smith,
and Medin (2012), and it has played an equally important role in 
artificial intelligence, especially the subfield of machine learning 
(Langley, 1996). Here is one candidate definition:

\vskip 0.05in

\bull
  {\it A concept is a cognitive structure that denotes a set of 
       entities or situations and that characterizes them in terms 
       of their regularities\/}.

\vskip 0.05in
\noindent
This statement is intentionally abstract, as needed if we want to cover 
  the variety of accounts that have been proposed. 
However, all treatments share the assumptions that concepts involve 
  mental structures, that these refer to some set or class, and that 
  this set is not arbitrary but rather is characterized by law-like
  descriptions. 
The terms {\it category\/} and {\it concept\/} may be used synonymously, 
  but the former often refers to the set being denoted and the latter 
  to mental structures that encode it. 

Humans associate words with many concepts in everyday life, most of them
  referring to classes of objects, whether inanimate (e.g., clouds, 
  street signs) or animate (e.g., dogs, cats), but not all of them are 
  linked to language. 
There can be considerable variation among the members of each category, 
  yet people group them together because they nevertheless share underlying 
  similarities. 
Specialists have their own concepts, from experts in automobile 
  diagnosis who distinguish types of faults to professional wine 
  tasters who discriminate among different vintages. 
Concepts help us encode, store, organize, and communicate our experience 
  of the world.   

% Representation - Rules, exemplars, probabilistic

If concepts are mental structures, then a cognitive system must
  {\it represent\/} them in some manner.
The literature reflects many competing positions on this issue, but 
  Medin and Coley (1998) have organized them into three broad 
  classes.\footnote{Medin and Coley refer to the first of these views 
  about conceptual representation as {\it classical\/} and the second 
  as {\it probabilistic\/}, but the labels {\it logical\/} and {\it 
  prototype\/} also appear in the literature.} 
These include: 

\vskip 0.05in
\noindent

\bull
{\it Logical accounts\/}, in which a concept specifies the necessary and 
  sufficient conditions for membership in a category.
This view, perhaps the earliest one, has roots in philosophy. 
However, treating concepts as logical conjunctions appears too 
  constraining for many human categories, and some variants allow 
  for disjunctive descriptions. 

\bull
{\it Prototype accounts\/}, in which a concept provides a generalized 
  description for a category that notes relevant features or relations 
  but allows for differing degrees of membership, with some instances 
  of a category fitting better than others. Probabilistic representations 
  of concepts are a widely adopted instance of this idea. 

\bull
{\it Exemplar accounts\/}, in which a concept comprises a set of 
  representative exemplars that, taken together, indicate membership 
  in the category. As with prototypes, some instances of a category 
  will fare better than other ones based on their similarity to the 
  stored exemplars, with notions of similarity differing across 
  alternative models. 

\vskip 0.05in
\noindent
Langley (1996) mentions a fourth option -- {\it threshold concepts\/} 
  -- that include some neural network accounts. 
There are variations on each theoretical position, often 
  introduced in response to observed phenomena in human behavior. 
There is not enough space to cover them all here, but the objective 
  of this paper is not to review them fully, only to 
  set the stage for a later comparison with chunks. 

% Organization - Set, tree, DAG

Some theories about concepts also make statements about the {\it 
  organization\/} of these conceptual structures in long-term memory. 
The simplest, sometimes implicit, position is that memory simply contains 
  a set of concepts, with no further constraints on their storage. 
Another account (e.g., Fisher, 1987) claims that concepts are organized 
  in a taxonomic hierarchy with more general categories above their 
  more specific descendants. 
A related but distinct alternative states that concepts serve as nodes in 
  a lattice or directed acyclic graph, which means that they can have 
  multiple parents. 
A third scheme organizes knowledge in an {\it inference network\/}
  (Langley, 1996), in which simpler concepts provide evidence for more 
  complex ones, with neural networks and logic programs being examples. 
These different frameworks provide competing indexing schemes for 
  conceptual structures. 

% Performance
% Matching, similarity, posteriors

As Anderson (1978) has argued, representational assumptions have little
  predictive or explanatory power without additional claims about how
  the structures are {\it used\/}. 
Thus, theories about conceptual knowledge typically come with postulates 
  about performance mechanisms that operate over them. 
In particular, accounts that state concepts are logic-like in character
  typically posit an all-or-none matching process to determine whether 
  they apply to given situations.
In contrast, prototype theories suppose some form of partial matching 
  that computes a similarity score or a probability for a concept 
  description given a 
  situation.\footnote{Note that theories of structural analogy (e.g., 
  Gentner \& Forbus, 2011) assume partial matching over logic-like 
  structures, which supports the argument that performance mechanisms 
  are crucial.} 
Finally, exemplar or case-based approaches apply the same basic idea 
  to calculate similarities between stored cases and a situation. 
Typical measures of performance are classification error (or accuracy)
  and recognition (response) time. 

% Learning

% Supervised / unsupervised, incremental / batch
% Generalization, discrimination, new exemplars, prob updates

There is also a substantial literature on concept {\it acquisition\/} 
in humans and machines. 
Work in both communities has concentrated on supervised training, in
  which each case has an associated class label, but some efforts 
  have addressed unsupervised learning (e.g., Fisher et~al., 1991). 
Research in machine learning has relied mainly on batch handling 
  of training cases, but it is clear that people process experiences 
  incrementally, so many models of human concept learning make this 
  assumption. 
A variety of incremental mechanisms have appeared in the literature, 
  including generalization from specific cases, specialization of 
  overly general structures, and probabilistic updates of statistical 
  descriptions. 
One characteristic of human acquisition that has received less
  attention is its piecemeal character (Langley, 2022), which 
  involves learning one structure at a time. 

% Literature on concepts (embed in earlier text?) 
% \note{*** Add paragraph on literature here or embed earlier?}

\section{Chunks in Cognitive Systems}

% Definition

The study of chunks has also received considerable attention in cognitive
  psychology (Miller, 1956; Fountain \& Doyle, 2012), although somewhat 
  less than concepts. 
The idea has also played a role in artificial intelligence, although 
  seldom by that name, with Campbell and Berliner's (1983) work on
  chunks in chess being an important exception. 
Here is a tentative definition: 

% starting with Miller (1956), 

\vskip 0.05in

\bull
  {\it A chunk is a cognitive structure that denotes a collection of 
       entities, possibly typed, that form a specified configuration 
       or relational pattern\/}. 

\vskip 0.05in
\noindent
Again, this statement is intentionally abstract in order to cover 
  the range of theories that have appeared in the literature. 
Nevertheless, despite their differences, these accounts share 
  assumptions that chunks are cognitive structures, they refer
  to collections of entities, and they are characterized as patterns 
  of relations among their constituents. 

Everyday examples of chunks include written words, telephone numbers, 
familiar tunes, standard table settings, and the layout of rooms in 
a house. They also arise in more specialized context, and some of the
strongest evidence for them comes from studies of expert memory. For 
instance, Chase and Simon (1973) found that chess masters store and use
chunks for many board configurations. Similarly, it is obvious that chefs
know the ingredients for many recipes, actors recall extended passages
from many plays, and experienced taxi drivers remember the layout for 
their home cities. Chunks are as pervasive as concepts in human cognition. 

% Representation

Consider the above definition's implications for mental representation. 
A chunk specifies a collection of entities that satisfy some relational
  pattern or configuration.
This does not constrain the constituent entities, so they may be physical 
  objects (e.g., silverware on a table), marks on a page (e.g., letters), 
  or even complex activities (e.g., dance steps). 
Neither does it limit the nature of the relations, which may be spatial
  (as in constellations), temporal (as in music), or something more 
  abstract. 
Moreover, these relations may be purely qualitative (e.g., above, behind) 
  or they may include quantitative details (e.g., distance, angle). 
Finally, the literature usually treats chunks as logic-like structures, 
  but nothing prevents them from including probabilities or other 
  annotations. 

% Organization

The definition also allows chunk constituents to have types, which means 
  the latter may themselves be other chunks. 
This has import for memory organization, as it implies that a chunk can 
  specify a hierarchical structure. 
However, this does not involve the taxonomic hierarchies that organize 
  some concepts, but rather a partonomic or compositional hierarchy.
For example, a phrase can be broken down into a sequence of words, each 
  of which comprises a sequence of letters. 
Not all treatments of chunks make this assumption, but they do not rule 
  it out and there is general agreement that high-level chunks may be
  specified in terms of lower-level ones. 

% \footnote{Some accounts organize chunks in a discrimination network
% (Richman, 1991), which has similarities to a taxonomy.} 

% Performance
% Recognition and reconstruction

Theories of chunk use in humans emphasize two primary capabilities. 
The first involves the {\it recognition\/} of chunks from perceptions, 
  such as a friend's telephone number, a familiar passage of music,
  or a famous constellation in the night sky. 
Accounts in the literature typically assume this is a bottom-up,
  all-or-none match process, with relations among elements playing 
  a key role. 
The second involves {\it recall\/} or reconstruction of chunks, typically 
  after some delay. 
Examples here include remembering a sequence of digits or words after 
  hearing them or recreating positions on a chess board after a brief 
  exposure. 
In this case, most accounts assume some form of top-down generative
  process. 
Typical measures here are memory capacity and reconstruction time. 

% Measures

% Add paragraph on measures of performance. 
% Memory capacity and reconstruction time. 

% Learning is incremental, unsupervised, cumulative.} 

There have been some computational accounts of chunk acquisition
  in humans, but, again, fewer than for concept learning. 
Computational models like EPAM (Richman, 1991) and CHREST (Gobet \& 
  Lane, 2012) assume, based on empirical evidence, that this process 
  is incremental and largely unsupervised, with mastering the ability 
  to recognize chunks preceding the ability to reconstruct them. 
Educational studies also suggest that it is typically cumulative, 
  with high-level chunks being acquired only after those for their 
  constituents have been learned.  
Research in machine learning on tasks like grammar induction 
  (e.g., Wolff, 1982; Langley \& Stromsten, 2000) has often used 
  nonincremental methods to introduce chunks, but these appear 
  to differ from those used by humans.\footnote{The `chunks' acquired 
  by the Soar architecture (Laird, 2012), which primarily encode 
  procedural knowledge, bear little resemblance to Miller's original 
  notion, which emphasized perception and memory.} 

\newpage

\section{Analysis of Similarities and Differences} 
% \vspace{-0.06in}

Now that we have reviewed key ideas from the literature on concepts and
chunks, we can compare and contrast them. As seen above, their study
has emphasized different aspects of cognition, but this does rule out 
the possibility that one type of mental structure, or even a single 
set of processes, underlies both sets of cognitive phenomena. Because
this is our main hypothesis, we should start by examining their common
assumptions. 

% Commonalities

We can identify five primary postulates that are shared by the majority
of psychological research on concepts and chunks. Many of these tenets
also hold for their treatment in artificial intelligence. They include:

\vskip 0.05in

\bull
{\it Discrete cognitive structures\/}. 
One common assumption is that concepts and chunks are both encoded
as discrete structures which are stored in long-term memory. A 
given concept or chunk may be linked to others, but they remain
separate mental entities with their own content. 

\bull
{\it Organization in memory\/}. 
Another shared postulate is that both concepts and chunks are
organized and indexed in long-term memory. The forms of their 
organization differ, but discrete structures of a given type
are nevertheless linked in some fashion. 

\bull
{\it Access through recognition\/}. 
Theories of concepts and chunks both posit they can be accessed
from long-term memory through a recognition process. This involves 
matching their descriptions against perceptions, but it may also 
take advantage of their organization. 

\bull
{\it Response time as a performance metric\/}. 
Empirical studies of concepts and chunks in humans often use
response time as a measure of performance. This involves
recognition time for concepts and reconstruction time for chunks,
but both equate rapid processing with mastery. 

\bull
{\it Incremental and piecemeal learning\/}. 
Studies of human learning reveal that it occurs incrementally, with 
one or a few instances being processed at a time, and in a piecemeal 
manner, with one or a few cognitive structures being acquired or 
revised at a time. 

% (Langley, 2022)

\vskip 0.05in
\noindent
These commonalities offer some encouragement that a unified account
of concepts and chunks is possible, and efforts toward this goal can 
use them as constraints on candidate theories, as already done in
many separate accounts. 

% Differences

But as we have seen, despite their underlying similarities, the
research community has generally treated concepts and chunks as 
though they are different types of mental entities. Some important
distinctions between them include: 

\vskip 0.05in

\bull
{\it Categories vs.\ composites\/}. 
Theoretical treatments of concepts stress their ability to represent 
sets of instances or categories, whereas accounts of chunks emphasize
their compositional character. Chunks may cover different instances,
but this is downplayed. Concepts may involve relations, but this is
not required, whereas relational patterns are central to chunks.

% Concepts => generality? 
% Chunks => specificity? 

\bull
{\it Taxonomic vs.\ partonomic hierarchies\/}. 
Not all theories of concepts discuss organization, but a common idea
is that they are positioned in a taxonomic hierarchy, with higher nodes 
being more general than their descendants. In contrast, theories of
chunks typically posit they are organized in a partonomic hierarchy,
with parents being composed of their children. 

\bull
{\it Flexible vs.\ strict matches\/}. 
Most accounts of conceptual processing assume flexible matching, with
classic examples being prototype and exemplar theories. Explanations
of chunk retrieval often treat their defining patterns as abstract 
yet rigid, with relations standing or falling together. 

% Concepts => Degrees of match
% Chunks => All-or-none match

\bull
{\it Recognition vs.\ reconstruction\/}.  
In most cases, studies of conceptual knowledge emphasize their
ability to {\it recognize\/} instances of categories, whereas
treatments of chunks usually emphasize their ability to 
{\it reconstruct\/} patterns from memory.\footnote{The distinction
between recognition and recall in cognitive psychology maps roughly 
onto the dichotomy between {\it discriminative\/} and {\it generative\/}
models in machine learning.} Some accounts of concepts discuss their
generative capacity, and chunks are certainly recognized, but these
tend to be downplayed. 

\bull
{\it Generalization vs.\ composition\/}.  
Treatments of concept acquisition emphasize generalization over
training cases. Learning occurs incrementally with each instance,
but multiple examples are often needed to establish category bounds.
Theories of chunk acquisition focus on composition of elements into
larger structures. Mastery may require multiple trials because elements
may be added incrementally. Moreover, chunking supports cumulative 
learning, in which later structures build on earlier ones, an idea
far less common with concepts. 

\vskip 0.05in
\noindent
Table~1 summarizes these differences, but it is important to note 
they are primarily distinctions in emphasis. Theories of conceptual
knowledge seldom forbid the features associated with chunks and
theories of chunks implicitly allow those associated with concepts.
This leaves the door open for an account that combines their features,
but we must still walk through this intellectual passage. 

\begin{table}[t]
\vskip -0.15in
\begin{center}
\caption{Different emphases in the treatments of concepts and chunks.}
\vskip 0.15in
\begin{tabular}{lllll}
\hline
\abovespace\belowspace
\quad Dimension      & ~~~~~~ & Concepts     & ~~~~~~ & Chunks \\
\hline
\abovespace
\quad Representation &        & Categories   &        & Composites \\
\quad Organization   &        & Taxonomic    &        & Partonomic \\
\quad Matching       &        & Flexible     &        & Strict \\
\quad Performance    &        & Recognition  &        & Reconstruction~~~ \\
\belowspace
\quad Acquisition    &        & Generalization &      & Composition \\
\hline
\end{tabular}
\end{center}
\end{table}

\section{A Proposed Research Agenda} 

% State goal of unifying two ideas, Explain why this is desirable 
% Note initial forays in this direction

Although a computational explanation of human-like intelligence that 
treats concepts and chunks separately may be possible, it would be
considerably less satisfying than a unified theory. There have 
been a few forays in the latter direction, but such attempts have 
received limited attention and the topic merits further effort. In
this section I consider some steps that the cognitive systems community
might take toward this end.

% One example is Richman's (1991) discussion of EPAM as a model of
% both concept formation and chunk acquisition. Another is Gobet and 
% Lane's (2012) report on the CHREST architecture, which extends a 
% theory of chunk acquisition to support concept-like `templates'. 

% Also discuss frames and scripts? 

\subsection{Representation} 

As noted, theories of concepts emphasize the representation of
generic categories or sets of items, whereas those for chunks
instead focus on their compositional character. A unified account 
would incorporate both ideas by positing a single type of cognitive
structure that describes a relational pattern of elements while also
characterizing a class of instances that satisfy it. Such structures
would have aspects of both concepts and chunks, although features of
one of the other would still be dominant in degenerate cases. 

For instance, a written word would be specified as a sequential
pattern of letters, but that word might have alternative spellings
which provide generality. Similarly, a structure for a chess board
would involve a spatial pattern of pieces that threaten or defend
each other, but different types of pieces might occupy the same role,
denoting a set of possible configurations. One could describe a passage
of music as a sequence of notes, but this could allow variations in
timing and emphasis in how it is played, again covering a range of
instances. 

% Explain how one might encode such generalized chunks. 

Of course, a full theory would need to specify the details of this
unified representation. In one version, some structures would be
more concept-like, with little compositional character, whereas
others would be more chunk-like, with little room for variation.
In another version, every structure would reflect both aspects, 
exhibiting both the generality of concepts and the configurational
nature of chunks. Whether these specifications are best encoded as
logical descriptions, prototypes, or exemplars remains an open question. 

\subsection{Organization}

We have seen researchers often assume that concepts are organized in a
taxonomic hierarchy, while chunks instead participate in partonomic
arrangements. A unified theory would incorporate both forms of mental
organization, with structures that refer to constituents and thus
impose a part-of hierarchy, but nevertheless reside within a taxonomy
in which higher-level nodes are generalizations of their descendants.
Human cognition supports both forms of organization and our computational
theory of intelligence should do so as well.

For example, the structure for mammals would include features like
hair covering and placentas, but it would also include constituents
like a head, torso, and limbs in a spatial configuration. Different
categories of mammals, such as quadrupeds and bipeds, would be
specializations that reside lower in the taxonomy. However, the 
latter would also contain nodes for categories of constituents, 
such as hoofed and clawed limbs, that appear in their concept-like
descriptions. Similarly, word classes would be specified as sequences 
of letters, with some variation, which in turn would be described by 
their arrangement, with greater variation due to handwriting style 
or printed fonts. 

% Explain how one might encode such generalized hierarchies.} 

A unified theory would specify how these two forms of organization 
interleave and complement one another. For instance, each node in
a taxonomic hierarchy could specify constituents that make up that 
concept, but these components would also reside in the same taxonomy 
or a kindred one. The account should also state whether the relations 
between chunk constituents are themselves decomposable or whether
they allow variation across instances. Formalisms like context-free
grammars and logic programs have some of the required features, but 
more appear necessary. 

% \subsection{Recognition and Recall}
\subsection{Performance} 

As noted above, many treatments of concept use favor flexible forms
of recognition, while accounts of chunk use emphasize recall or 
reconstruction. A unified theory would incorporate both capabilities,
relying on a form of partial matching when comparing stored structures
to perceptions and drawing on their relational patterns for reconstruction 
from long-term memory. Most theories of concepts and chunks do not 
rule this out, but their focus on different phenomena means they have
each told only part of the overall story. 

For instance, a bottom-up recognition process could group strings
of letters into words and then words into phrases to construct parse
trees for sentences, producing the same hierarchical structures for 
different fonts and spellings. In addition, a recall process could 
fill in missing letters or generate words for incomplete sentences.
The same mechanisms could recognize familiar configurations in chess
boards, even if they involve novel pieces, and reconstruct their layout
on the board from a chunk identifier and its constituent entities. 

In a unified account, structures would be recognized by finding their
degree of match to perceptions, possibly by sorting them through a
taxonomic hierarchy that indexes them. This process might occur in a
bottom-up manner, with categorization of constituents like letters
leading in turn to recognition of composites like words, but top-down
forms of retrieval might also be possible. Reconstruction would involve
the decomposition of structures into constituents, possibly through
multiple levels, using stored relations to organize the generative
process. 

\subsection{Acquisition}

% Two interleaved learning mechanisms? 

We have seen that theories of concept learning emphasize generalization
over experience, whereas accounts of chunking focus on creation of new
relational structures composed of elements. A unified framework would
incorporate both forms of acquisition, letting it create generalized
descriptions that cover distinct training cases but that also specify
relations among their constituents. This need not rely on a single
learning mechanism, but if it involves multiple processes, they must
work in tandem to support both facets of acquisition. 

For example, this cooperative arrangement would support learning of
letter `concepts', which can vary substantially in appearance, but
also construction of word `chunks' with letters as components.
However, different words might occur in similar contexts, leading 
to `concepts' like nouns and adjectives, which can then be
combined into `chunks' for adjectival phrases and noun phrases.
Such conjoined processes would support cumulative learning in which
lower levels of the partonomic hierarchy (each embedded in a 
taxonomy) would provide elements to compose higher levels. 

% Explain how such learning might operate (e.g., as in GRIDS).

A unified theory of concepts and chunks would explain how to acquire
these two-faceted structures. A crucial question is whether two
learning mechanisms are necessary, as in research on the induction of
context-free grammars (e.g., Wolff, 1982; Langley \& Stromsten, 2000), 
or whether a single process can suffice. The extended framework must 
acquire both concept-like and chunk-like features of new structures, 
but this might result from interactions between one learning process 
and multiple performance mechanisms.

\section{Three Candidate Theories}

We can explore the prospects for a unified theory by examining existing
frameworks that appear to hold potential for handling both concepts
and chunks. In this section, I review three candidate's assumptions
about representation, organization, performance, and learning, then
discuss how researchers might extend their coverage to both sets of
phenomena. 

\subsection{Discrimination Networks}

One of the earliest computational models of human expertise was EPAM 
(Feigenbaum, 1961; Feigenbaum \& Simon, 1984), which organized knowledge 
in a discrimination network. This takes the form of a tree, with each 
node indicating a test or attribute and each downward branch specifying 
an outcome or value, much like a decision tree (Quinlan, 1986). Each 
terminal node stores not a class label but an `image' or description 
(e.g., a conjunction of attribute values). Performance involves sorting 
a stimulus downward through the network until reaching a terminal
node and using its image to recall unspecified values. Learning is 
interleaved with performance, with a {\it discrimination\/} process 
adding a new branch when the stimulus matches no branch's value and
a {\it familiarization\/} mechanism adding details to an image on 
reaching a terminal node.

The earliest versions of EPAM focused on rote memorization, as required
by tasks like paired-associates learning. However, later ones (e.g.,
Richman, 1991) augmented the framework to let it store images at internal 
nodes, making their discrimination networks similar to taxonomies of 
concepts. CHREST (Gobet \& Lane, 2012; Lane, Gobet, \& Smith, 2008)
further extended these images into `templates' whose roles can be 
filled in different ways, making them even more like generalized
concepts. Both Richman and Lane et~al.\ claimed that their discrimination
networks supported chunk storage, use, and acquisition, but they seem 
to have equated chunks with images, and did not provide a way to acquire
high-level chunks in terms of simpler ones.

We should consider what must be added to support the compositional
character of chunks and their cumulative acquisition. EPAM's tests
were not limited to primitive features; they could refer to nodes
elsewhere in the network. This meant that it could define complex chunks
(e.g., words as sequences of letters) in terms of simpler ones (e.g.,
letters as configurations of lines). However, in such cases, EPAM 
and its successors appear to have required stimuli that were already
organized into partonomic structures. What they lacked was some means
to acquire these structures from training data and some way to infer
them during performance. Techniques from grammar induction (e.g., Wolff,
1982) that introduce nonterminal symbols offer one response to the 
first issue and classic methods for parsing offer an approach to the 
second. Integrating these with the framework of discrimination networks 
could offer a unified account of concepts and chunks.

\vspace{-0.04in}
\subsection{Neural Networks}

A second promising paradigm, multilayer neural networks, organizes 
expertise in a very different manner. Nodes denote features or
attributes and weighted links connect them to others, typically
in a directed acyclic graph. Input nodes at the lowest level feed
into intermediate ones, often in many layers, leading ultimately
to output nodes at the highest level. Performance involves passing
values of perceived features to input nodes, using the weighted links
to compute values for intermediates, and finally calculating values 
for output nodes. For classification tasks, the system typically
selects the highest-scoring node to assign a class. Learning usually
relies on some variant of backpropagation (Rumelhart et~al., 1986),
which alters weights on links to reduce errors in output values.

Because neural networks are widely used for classification tasks,
they have clear relevance to concepts, and Kruschke's (1992) ALCOVE
was an early model of human categorization cast in these terms.
However, we can also view their internal nodes as encoding chunk-like
structures, although this is not readily apparent because they take on
continuous values, whereas classic chunks are discrete. Yet on visual 
processing tasks like letter and face recognition, examination of the
weights on hidden units show they correspond to constituent features 
for edges, corners, and other facets of images. This suggests that
the widespread claim neural networks exhibit {\it representation
learning\/} by creating latent variables is equivalent to stating
they acquire chunks. 

Based on this observation, we might conclude that neural networks
already provide a unified account of concepts and chunks, even
though this idea has not received attention in the literature.
The fact that they must be initialized with a complete network
rather than adding nodes (chunks) as needed is not an obstacle, as 
starting weights on links to intermediate nodes can be negligible.
The problem lies not with representation or performance, but with 
learning, as humans acquire chunks and concepts far more rapidly
than the gradient descent methods that neural networks use to
update their parameters. To support human-like behavior, the
framework will need alternative mechanisms that learn in a more
rapid and piecemeal manner (Langley, 2022). 

\vspace{-0.05in}
\subsection{Probabilistic Concept Hierarchies}
\vspace{-0.01in}

A third theoretical framework, first introduced in Cobweb (Fisher,
1987; Iba \& Langley, 2011), focused explicitly on concept 
formation. This assumes that concepts are organized as nodes in a 
taxonomic hierarchy, with terminal nodes denoting training cases
and nonterminals describing probabilistic summaries of their
descendants. A performance mechanism sorts new cases down the path
that maximizes an information-theoretic metric, predicting missing
features as needed. Learning is incremental and unsupervised,
with each node's distribution being updated as instances pass 
through it and new concepts being created when reaching a terminal 
node or lacking a good-scoring child. The framework combines ideas 
from discrimination network, probabilistic, and exemplar accounts, 
and it exhibits behavior consistent with some categorization phenomena.

The variants of Cobweb developed to date have focused on conceptual
knowledge, but some extensions could let the framework support chunk
representation, organization, use, and acquisition as well. To this
end, they might assume that instances comprise a set of elements
(e.g., words or objects), each with a description but also linked
by relations (e.g., order or direction), and that concepts specify 
a set of constituents and their surrounding contexts, along with the
probability distributions for each. They might also include a parsing
mechanism that, when processing an instance, creates candidate chunks
that it sorts through the taxonomy to select the best, which it then
uses to elaborate the instance. Every nonprimitive concept would have
constituents and thus also be a chunk. Parsing would use this knowledge
to recode perceived instances by adding partonomic structure. Learning
would remain largely the same except for operating over the extended 
representation.

% A learning mechanism that incorporates selected chunks, including
% their constituents and contexts, into the taxonomy, updating the 
% probabilities and altering organization as needed.

The commitment that all but primitive training cases contain multiple
elements would ensure the constituents that provide material for 
parsing and chunk creation. The inclusion of context would allow 
creation of concepts based on this information rather than ones based
only on similar content. This would let Cobweb acquire grammatical
knowledge by identifying words (black, cat) as letter
sequences, word classes (adjective, noun) in terms of shared contexts,
phrasal structures as sequences of word classes, and phrasal classes
that arise in shared contexts at a higher level. Presumably, the same
mechanisms could explain chunk processing in domains like chess.
Chunks in Cobweb would include probabilities, but highly regular
ones would behave like logical structures, as in classic accounts.
The resulting theory would unify concepts and chunks. 

% Some approaches to the induction of context-free grammars (e.g.,
% Wolff, 1982) alternate between the creation of nonterminal symbols
% (chunks) based on recurring content (words or phrases) and then 
% merging these symbols them based on similar context (word or phrasal
% classes). However, such systems make no contact with the literature
% on categorization and concept learning, and they are typically
% nonincremental. The design for an extended Cobweb offers one avenue,
% although not the only one, for unifying these two distinct traditions.

\section{Summary Remarks}

In this paper, I noted the cognitive system paradigm's emphasis
on cognitive structures and examined two varieties -- concepts and
chunks -- that have received attention in cognitive psychology.
I reviewed common assumptions about the representation, organization,
use, and acquisition of concepts and chunks, then examined similarities
and differences in their treatments. I concluded that these two types 
of entities may not be not truly distinct, but rather simply facets 
of the same mental encodings viewed from different perspectives. 

In response, I proposed a research agenda for developing unified 
theories of concepts and chunks. This would explore generalized 
representations that subsume both ideas, common organizations 
for them in long-term memory, unified performance mechanisms that 
operate over them, and learning processes that acquire and refine 
them. I also discussed three existing frameworks that seem ripe
for extension. We should evaluate each in terms of its ability to
support abilities associated with both concepts and chunks, subject 
to constraints on human cognition. Progress in this direction would
take us an important step closer to a unified theory of the mind.

\vspace{-0.04in}
\section*{Acknowledgements}
\vspace{-0.01in}

Work on this essay was supported by Grant FA9550-23-1-0580 from the
US Air Force Office of Scientific Research and by Grant N00014-23-1-2525
from the Office of Naval Research, which are not responsible for its
contents. I thank Chris MacLellan, Doug Fisher, and many others for
discussions that contributed to the ideas presented here. 

\vskip 0.20in
\noindent
\leftline{\large\bf References}
\vskip 0.10in

{\parindent -10pt\leftskip 10pt

Anderson, J.\ R.\ (1978). Arguments concerning representations for
mental imagery. {\it Psychological Review\/}, {\it 85\/}, 249--277.

Campbell, M., \& Berliner, H.\ J.\ (1983). A chess program that chunks.
{\it Proceedings of the Third National Conference on Artificial 
Intelligence\/} (pp.\ 49--53). Washington, DC: AAAI Press. 

Chase, W.\ G., \& Simon, H.\ A.\ (1973). Perception in chess. {\it
Cognitive Psychology\/}, {\it 4\/}, 55--81.

Feigenbaum, E.\ A.\ (1961). The simulation of verbal learning behavior. 
{\it Proceedings of the Western Joint Computer Conference\/} (pp.\ 
121--132). Los Angeles, CA. 

Feigenbaum, E.\ A., \& Simon, H.\ A.\ (1984). EPAM-like models of 
recognition and learning. {\it Cognitive Science\/}, {\it 8\/},
305--336. 

Fisher, D.\ H.\ (1987). Knowledge acquisition via incremental conceptual
clustering. {\it Machine Learning\/}, {\it 2\/}, 139--172. 

Fisher, D.\ H., Pazzani, M.\ J., \& Langley, P. (Eds.) (1991). 
{\it Concept formation: Knowledge and experience in unsupervised
learning\/}. San Francisco, CA: Morgan Kaufmann.

Fountain, S.\ B., \& Doyle, K.\ E.\ (2012). Learning by chunking.
In N.\ M.\ Seel (Ed.), {\it Encyclopedia of the sciences of learning\/}. 
Springer, Boston, MA.

Gentner, D., \& Forbus, K.\ D.\ (2011). Computational models of
analogy. {\it Wiley Interdisciplinary Reviews: Cognitive Science\/},
{\it 2\/}, 266--276.

Gobet, F., \& Lane, P.\ C.\ R.\ (2012). Chunking mechanisms and 
learning. In N.\ M.\ Seel (Ed.), {\it Encyclopedia of the sciences 
of learning\/}. New York, NY: Springer. 

Iba, W.\ F., \& Langley, P.\ (2011). Cobweb models of categorization 
and probabilistic concept formation. In E.\ M.\ Pothos \& A.\ J.\ 
Willis (Eds.), {\it Formal approaches in categorization\/}. Cambridge:
Cambridge University Press.

Kruschke, J.\ K.\ (1992). ALCOVE: An exemplar-based connectionist
model of category learning. {\it Psychological Review\/}, {\it 99\/},
22--44.

Laird, J.\ E.\ (2012). {\it The Soar cognitive architecture\/}. 
Cambridge, MA: MIT Press.

Lane, P.\ C.\ R., Gobet, F., \& Smith, R.\ L.\ (2008). Attention 
mechanisms in the CHREST cognitive architecture. {\it Proceedings 
of the Fifth International Workshop on Attention in Cognitive
Systems\/} (pp.\ 207--220). Graz, Austria: Joanneum Research.

Langley, P.\ (2012). The cognitive systems paradigm. {\it Advances in
Cognitive Systems\/}, {\it 1\/}, 3--13.

Langley, P.\ (2022). The computational gauntlet of human-like learning.
{\it Proceedings of the Thirty-Sixth AAAI Conference on Artificial 
Intelligence\/} (pp.\ 12268--12273). Vancouver: AAAI Press.

% Vancouver, BC: 

Langley, P., Laird, J.\ E., \& Rogers, S.\ (2009). Cognitive
architectures: Research issues and challenges. {\it Cognitive 
Systems Research\/}, {\it 10\/}, 141--160.

Langley, P., \& Stromsten, S.\ (2000). Learning context-free grammars 
with a simplicity bias. {\it Proceedings of the Eleventh European 
Conference on Machine Learning\/} (pp.\ 220--228). Barcelona, Spain: 
Springer-Verlag.

% Medin, D.\ L., \& Rips, L.\ J.\ (2005). Concepts and categories:
% Memory, meaning, and metaphysics. In K.\ J.\ Holyoak \& R.\ G.\ 
% Morrison (Eds.), {\it The Cambridge handbook of thinking and 
% reasoning\/} (pp. 37--72). Cambridge University Press.

Medin, D. L., \& Coley, J. D. (1998). Concepts and categorization.
In J.\ Hochberg (Ed.), {\it Perception and cognition at century's
end\/}, 403--439. New York, NY: Academic Press.

Miller, G.\ A.\ (1956). The magical number seven, plus or minus
two: Some limits on our capacity for processing information.
{\it Psychological Review\/}, {\it 63\/}, 81--97.

Quinlan, J.\ R.\ (1986). Induction of decision trees. {\it Machine 
Learning\/}, {\it 1\/}, 81--106.

Richman, H.\ B.\ (1991). Discrimination net models of concept 
formation. {\it Concept formation: Knowledge and experience in 
unsupervised learning\/}. San Francisco, CA: Morgan Kaufmann.

Rips, L.\ J., Smith, E.\ E., \& Medin, D.\ L. (2012). Concepts and
categories: Memory, meaning, and metaphysics. In K.\ J.\ Holyoak \&
R.\ G.\ Morrison (Eds.), {\it The Oxford handbook of thinking and
reasoning\/}, 177--209. Oxford, UK: Oxford University Press. 

Rumelhart, D.\ E., Hinton, G.\ E., \& Williams, R.\ J.\ (1986). Learning
internal representations by back-propagating errors. In D.\ E.\ Rumelhart
\& J.\ L.\ McClelland (Eds.), {\it Parallel distributed processing\/}
(Vol. 1), 318-362. Cambridge, MA: MIT Press.

% Smith, E.\ E., \& Medin, D.\ L.\ (1981). {\it Categories and
% concepts\/}. Cambridge, MA: Harvard University Press.

Wolff, J.\ G.\ (1982). Language acquisition, data compression 
and generalization. {\it Language \& Communication\/}, {\it 2\/}, 
57--89. 

}

\end{document}

